\section{Introdução}
\label{sec:introducao}

O abastecimento é uma atividade cotidiana, mas envolve uma decisão que precisa
ser correta antes do início da transferência do combustível: associar o produto
selecionado à configuração admitida pelo veículo. Quando essa associação falha,
o procedimento posterior depende do tipo e da quantidade de combustível, do
funcionamento do motor e das recomendações do fabricante. A literatura registra
que o erro não pode ser explicado apenas por falta de atenção. Em uma análise
de tarefa e entrevistas com 23 pessoas que haviam abastecido veículos leves com
o produto errado no Reino Unido, \cite{adamsDavid2007} identificaram combinações
de fatores humanos e ambientais capazes de formar caminhos para o erro. Os
autores defenderam intervenções no leiaute e no equipamento para interromper
esses caminhos.

O presente trabalho nasce de um problema relatado em uma rede de postos de
combustíveis do sul do Brasil: durante o atendimento, pode ocorrer a seleção de
um produto incompatível com o veículo. A existência, frequência e composição
dos custos locais ainda serão verificadas por registros administrativos; por
isso, esta versão não apresenta número de ocorrências nem valor de prejuízo.
Essa separação é importante, pois dados internacionais ou relatos isolados não
podem ser convertidos diretamente em prevalência para a organização estudada.

Como possível barreira anterior à liberação, propõe-se consultar a placa do
veículo em um serviço autorizado, transformar o código retornado em uma política
explícita de compatibilidade e liberar temporariamente apenas o bico solicitado
e compatível. O serviço Consulta Online SENATRAN disponibiliza código e descrição
do combustível na modalidade detalhada \cite{serproDadosVeiculo2026}. Entretanto,
seu uso não se resume a uma chamada técnica: requer caso de uso aprovado,
contratação do Serpro e de uma Gerenciadora de Consentimento e Ciência (GCC) e
atendimento às condições de governança aplicáveis \cite{senatran1392025,serproComoContratar2026}.

O artefato acadêmico será inicialmente construído em bancada seca, sem
combustível e em baixa tensão. A trava física representa o intertravamento, mas
não constitui solução pronta para instalação em uma bomba. Uma implantação real
exigiria integração com o fabricante ou controlador de pista, classificação do
local e equipamentos adequados à atmosfera potencialmente explosiva
\cite{inmetro1152022,nr20}.

\subsection{Problema e pergunta de pesquisa}

O problema de pesquisa consiste na ausência de uma barreira automatizada que,
no contexto estudado, verifique a compatibilidade antes de autorizar o acesso ao
bico. A investigação é orientada pela seguinte pergunta:

\begin{quote}
\textit{Como um sistema baseado em consulta veicular e intertravamento pode
impedir a liberação de combustíveis incompatíveis em ambiente controlado, e em
quais condições sua adoção apresenta viabilidade econômica para uma rede de
postos?}
\end{quote}

Foram formuladas duas hipóteses verificáveis:

\begin{itemize}
  \item \textbf{H1:} na matriz de testes de bancada, o sistema impedirá todas as
  liberações incompatíveis e não liberará automaticamente nenhum bico diante de
  resposta inválida, ambígua ou indisponível;
  \item \textbf{H2:} em pelo menos um dos cenários de sensibilidade, as perdas
  evitáveis estimadas superarão o custo anualizado da solução.
\end{itemize}

H1 somente será aceita se não houver liberação indevida nos casos definidos
antes dos ensaios. H2 será avaliada por cenários, e não como economia observada.
Ambas permanecem sem teste conclusivo nesta etapa.

\subsection{Objetivos}

O objetivo geral é projetar, implementar e avaliar um protótipo para prevenção
de abastecimento incompatível, integrando consulta veicular, controle de
autorização e análise econômica baseada em dados históricos anonimizados.

São objetivos específicos:

\begin{enumerate}
  \item caracterizar o processo de abastecimento e, quando existirem registros
  suficientes, as ocorrências e os custos documentados no estudo de caso;
  \item mapear evidências científicas e documentação técnica sobre erro de
  abastecimento, fatores humanos, identificação veicular e intertravamento;
  \item elicitar requisitos funcionais, não funcionais, de privacidade e de
  segurança para o artefato;
  \item desenvolver um simulador do provedor veicular e um protótipo em bancada
  com política de falha para o estado bloqueado;
  \item avaliar o comportamento funcional, o tratamento de falhas, a latência,
  a usabilidade e a repetibilidade do protótipo; e
  \item estimar custos e ponto de equilíbrio sob diferentes percentuais de
  prevenção, distinguindo gastos comprovados de premissas.
\end{enumerate}

\subsection{Justificativa e contribuições esperadas}

A relevância prática decorre do potencial de agir antes do dano, quando a
escolha ainda pode ser bloqueada. A relevância acadêmica está na construção e
avaliação explícita de um artefato sociotécnico: não basta reconhecer a placa;
é necessário tratar códigos desconhecidos, indisponibilidade, exceções
operacionais, dados pessoais e o limite entre protótipo e equipamento de campo.

As contribuições esperadas são: (i) uma caracterização documentada do problema
na empresa; (ii) um protocolo e uma síntese da literatura; (iii) requisitos e
arquitetura auditáveis; (iv) um protótipo reprodutível de bancada; (v) uma matriz
de testes com estados de falha; e (vi) um modelo econômico transparente. Essas
contribuições são metas do trabalho e não resultados já alcançados.
